% ==========================================================
%  ML UNSUPERVISED TECHNIQUES — MASTER SUMMARY
%  Based on: Hochreiter, Machine Learning: Unsupervised Methods, JKU Linz
% ==========================================================

\documentclass[11pt, oneside]{book}

\usepackage[
  a4paper,
  top    = 2.6cm,
  bottom = 2.1cm,
  left   = 2.5cm,
  right  = 2.5cm,
  headheight = 15pt
]{geometry}

\usepackage[T1]{fontenc}
\usepackage[utf8]{inputenc}
\usepackage{lmodern}
\usepackage{microtype}
\usepackage{inconsolata}

\setlength{\parindent}{0pt}
\setlength{\parskip}{0.62em}

\usepackage{amsmath, amssymb, amsthm, mathtools}
\usepackage{bm}
\usepackage{bbm}
\usepackage{mathrsfs}
\usepackage{nicefrac}
\usepackage{cancel}

\usepackage{siunitx}
\sisetup{per-mode = symbol}

\usepackage{graphicx}
\usepackage{booktabs}
\usepackage{array}
\usepackage{float}
\usepackage{enumitem}

\setlist[itemize]{topsep=3pt, itemsep=2pt, leftmargin=2.2em}
\setlist[enumerate]{topsep=3pt, itemsep=2pt, leftmargin=2.4em}

\usepackage[ruled, vlined, linesnumbered]{algorithm2e}
\SetAlFnt{\small}
\SetAlCapFnt{\small\bfseries}
\SetCommentSty{textit}

\usepackage{xcolor}

\definecolor{accent}{HTML}{1B3A5C}
\definecolor{def@frame}{HTML}{1A5276}   \definecolor{def@bg}{HTML}{EBF5FB}
\definecolor{thm@frame}{HTML}{922B21}   \definecolor{thm@bg}{HTML}{FDEDEC}
\definecolor{ex@frame} {HTML}{1E8449}   \definecolor{ex@bg} {HTML}{EAFAF1}
\definecolor{int@frame}{HTML}{9A6A00}   \definecolor{int@bg}{HTML}{FEF9E7}
\definecolor{note@frame}{HTML}{5B2C6F}  \definecolor{note@bg}{HTML}{F5EEF8}

\usepackage{listings}
\lstset{
  basicstyle       = \ttfamily\small,
  breaklines       = true,
  numbers          = left,
  numberstyle      = \tiny\color{gray!70},
  frame            = tb,
  framerule        = 0.4pt,
  rulecolor        = \color{gray!35},
  tabsize          = 2,
  showstringspaces = false,
  keywordstyle     = \color{def@frame}\bfseries,
  commentstyle     = \color{gray!55!black}\itshape,
  stringstyle      = \color{ex@frame!90!black},
  backgroundcolor  = \color{gray!5},
  xleftmargin      = 2em,
  xrightmargin     = 0.4em,
  aboveskip        = 0.8em,
  belowskip        = 0.5em
}

\usepackage[most]{tcolorbox}

\tcbset{
  enhanced,
  breakable,
  arc             = 2.5mm,
  boxrule         = 0.5pt,
  left            = 4.5mm,
  right           = 4mm,
  top             = 2.5mm,
  bottom          = 2.5mm,
  fonttitle       = \bfseries\small,
  sharp corners   = west,
  parbox          = false
}

\newtcolorbox{defbox}[1][]{
  title  = {Definition},
  colback = def@bg, colframe = def@frame,
  borderline west = {3.5pt}{0pt}{def@frame},
  #1
}

\newtcolorbox{thmbox}[1][]{
  title  = {Theorem / Key Formula},
  colback = thm@bg, colframe = thm@frame,
  borderline west = {3.5pt}{0pt}{thm@frame},
  #1
}

\newtcolorbox{exbox}[1][]{
  title  = {Example},
  colback = ex@bg, colframe = ex@frame,
  borderline west = {3.5pt}{0pt}{ex@frame},
  #1
}

\newtcolorbox{intbox}[1][]{
  title  = {Intuition \& Mental Model},
  colback = int@bg, colframe = int@frame,
  borderline west = {3.5pt}{0pt}{int@frame},
  #1
}

\newtcolorbox{notebox}[1][]{
  title  = {Note},
  colback = note@bg, colframe = note@frame,
  borderline west = {3.5pt}{0pt}{note@frame},
  #1
}

\usepackage[titles]{tocloft}

\renewcommand{\cftchapfont}     {\bfseries\color{accent}}
\renewcommand{\cftchappagefont} {\bfseries\color{accent}}
\renewcommand{\cftchapleader}   {\bfseries\color{accent!40}\cftdotfill{\cftdotsep}}
\setlength{\cftbeforechapskip} {7pt}

\renewcommand{\cftsecfont}      {\small}
\renewcommand{\cftsecpagefont}  {\small}
\setlength{\cftsecindent}      {1.8em}
\setlength{\cftbeforesecskip}  {1.5pt}

\renewcommand{\cftsubsecfont}      {\small\color{gray!65!black}}
\renewcommand{\cftsubsecpagefont}  {\small\color{gray!65!black}}
\setlength{\cftsubsecindent}      {3.4em}
\setlength{\cftbeforesubsecskip}  {0.5pt}

\usepackage{titlesec}

\titleformat{\chapter}[display]
  {\normalfont\bfseries\color{accent}\huge}
  {\hfill\fontsize{68}{68}\selectfont\color{accent!13}\thechapter}
  {-10pt}
  {}
  [{\vspace{4pt}%
    {\color{accent}\rule{\linewidth}{1.6pt}}\\[-2pt]%
    {\color{accent!35}\rule{\linewidth}{0.5pt}}}]
\titlespacing*{\chapter}{0pt}{-28pt}{22pt}

\titleformat{\section}
  {\normalfont\large\bfseries\color{accent}}
  {\thesection}
  {0.8em}
  {}
  [{\vspace{1pt}\color{accent!45}\titlerule[0.45pt]}]
\titlespacing*{\section}{0pt}{3.8ex plus 1ex minus .2ex}{2.2ex plus .2ex}

\titleformat{\subsection}
  {\normalfont\normalsize\bfseries\color{accent!80!black}}
  {\thesubsection}
  {0.7em}
  {}
\titlespacing*{\subsection}{0pt}{2.6ex plus 1ex minus .2ex}{1.4ex plus .2ex}

\titleformat{\paragraph}[runin]
  {\normalfont\small\bfseries}
  {}
  {0em}
  {}[.\enspace]
\titlespacing{\paragraph}{0pt}{1.2ex plus .5ex minus .2ex}{0.5em}

\usepackage{fancyhdr}
\pagestyle{fancy}
\fancyhf{}
\fancyhead[L]{\small\color{gray!70}\nouppercase{\leftmark}}
\fancyhead[R]{\small\color{gray!70}\thepage}
\renewcommand{\headrulewidth}{0.4pt}

\fancypagestyle{plain}{%
  \fancyhf{}
  \fancyfoot[C]{\small\color{gray!70}\thepage}
  \renewcommand{\headrulewidth}{0pt}%
}

\usepackage[hidelinks]{hyperref}
\usepackage[nameinlink, capitalize]{cleveref}


% ══════════════════════════════════════════════════════════
%  DOCUMENT METADATA
% ══════════════════════════════════════════════════════════
\newcommand{\coursetitle}{ML Unsupervised Techniques}
\newcommand{\documenttype}{Master Summary}
\newcommand{\authorname}{Stefan Eder}
\newcommand{\basedon}{Based on Hochreiter, \emph{Machine Learning: Unsupervised Methods}, JKU Linz}


% ══════════════════════════════════════════════════════════
%  MATHEMATICAL MACROS
% ══════════════════════════════════════════════════════════

% ── Number Sets ───────────────────────────────────────────
\newcommand{\R}{\mathbb{R}}
\newcommand{\N}{\mathbb{N}}
\newcommand{\Z}{\mathbb{Z}}
\newcommand{\C}{\mathbb{C}}

% ── Probability & Statistics ──────────────────────────────
\newcommand{\E}{\mathbb{E}}
\newcommand{\Prob}{\mathbb{P}}
\newcommand{\Var}{\mathrm{Var}}
\newcommand{\Cov}{\mathrm{Cov}}
\newcommand{\iid}{\overset{\mathrm{iid}}{\sim}}
\newcommand{\ind}{\mathbbm{1}}

% ── Information Theory ────────────────────────────────────
\newcommand{\KL}[2]{D_{\mathrm{KL}}\!\left(#1 \,\|\, #2\right)}
\newcommand{\Ent}{\mathrm{H}}
\newcommand{\MI}{\mathrm{I}}

% ── Calculus & Optimization ───────────────────────────────
\newcommand{\argmax}{\operatorname*{arg\,max}}
\newcommand{\argmin}{\operatorname*{arg\,min}}
\newcommand{\grad}{\nabla}
\newcommand{\dd}{\,\mathrm{d}}
\newcommand{\pd}[2]{\frac{\partial #1}{\partial #2}}
\newcommand{\pdd}[2]{\frac{\partial^2 #1}{\partial #2^2}}

% ── Linear Algebra ────────────────────────────────────────
\newcommand{\T}{^{\top}}
\newcommand{\inv}{^{-1}}
\newcommand{\tr}{\mathrm{tr}}
\newcommand{\diag}{\mathrm{diag}}
\newcommand{\norm}[1]{\left\lVert #1 \right\rVert}
\newcommand{\abs}[1]{\left\lvert #1 \right\rvert}
\newcommand{\inner}[2]{\left\langle #1,\, #2 \right\rangle}

% ── Course-specific notation ──────────────────────────────
% Parameter vector and estimate
\newcommand{\wvec}{\bm{w}}
\newcommand{\what}{\hat{\bm{w}}}
% Data vector and matrix
\newcommand{\xvec}{\bm{x}}
\newcommand{\Xmat}{\bm{X}}
% Latent / source variable
\newcommand{\uvec}{\bm{u}}
% Likelihood and free energy
\newcommand{\lhood}{\mathcal{L}}
\newcommand{\Fbound}{\mathcal{F}}
% MSE and bias operators
\newcommand{\mse}{\mathrm{mse}}
\newcommand{\bias}{b}
% Distributions
\newcommand{\Normal}{\mathcal{N}}

% ── Misc ─────────────────────────────────────────────────
\newcommand{\given}{\,\vert\,}
\newcommand{\st}{\;\text{s.t.}\;}
\newcommand{\eqdef}{\coloneqq}


% ==========================================================
%  BEGIN DOCUMENT
% ==========================================================
\begin{document}

\frontmatter


% ── Title Page ────────────────────────────────────────────
\thispagestyle{empty}
\vspace*{1.8cm}

\begin{center}
  {\color{accent}\rule{\linewidth}{1.8pt}}
  \vspace{0.65cm}

  {\fontsize{28}{34}\selectfont\bfseries\color{accent}\coursetitle\par}
  \vspace{0.25cm}
  {\large\color{gray!80!black}\documenttype\par}

  \vspace{0.55cm}
  {\color{accent!35}\rule{\linewidth}{0.5pt}}
  \vspace{0.55cm}

  {\large\authorname\par}
  \vspace{0.2cm}
  {\small\color{gray!70}\basedon\par}
\end{center}

\vspace{1.6cm}

\begin{center}
\begin{minipage}{0.84\textwidth}
  \small\color{gray!80!black}
  This document is a compact summary of the course.
  It aims to cover all relevant concepts and results, but long derivations,
  detailed proofs, and worked examples are often condensed or omitted in
  favour of clarity and conciseness.
  Use it as a study aid alongside the lecture notes and slides, not as a
  replacement for them.

  \medskip
  If you spot an error or a gap, please open an issue or pull request at
  \textbf{github.com/stevetgit} --- corrections are always welcome.

  \medskip
  \color{gray!60!black}
  Good study materials should be shared, not gatekept.
  If this helped you, take it as a reason to help the next colleague you see struggling with something.
\end{minipage}
\end{center}

\vfill
{\centering\small\color{gray!65} Compiled \today\par}
\clearpage


% ── Table of Contents ─────────────────────────────────────
\setcounter{tocdepth}{2}

\begingroup
  \let\cleardoublepage\relax
  \let\clearpage\relax
  \tableofcontents
\endgroup

\mainmatter


% ==========================================================
\chapter{Introduction}
% ==========================================================

\section{Machine Learning Introduction}

Machine learning is currently a major research topic in academia and industry alike.
Applications span computer vision, speech recognition, recommender systems, and the
analysis of large biological datasets.
In the biomedical context, modern measurement technologies --- mRNA concentrations
measured via microarrays or next-generation sequencing, SNP arrays, copy number
variation data --- create a demand for methods that can extract structure from
high-dimensional data without relying on labeled training examples.

The machine learning course series at JKU Linz comprises four complementary courses:
\begin{itemize}
  \item \textbf{Basic Methods of Data Analysis:} summary statistics, PCA, ICA,
        multidimensional scaling, clustering, biclustering.
  \item \textbf{ML Supervised Methods:} classification and regression, SVMs, neural
        networks, deep learning, LSTM, ARMA/ARIMA.
  \item \textbf{ML Unsupervised Methods} (this course): MLE, MAP, maximum entropy,
        EM, PCA, ICA, factor analysis, mixture models, HMMs, Boltzmann machines.
  \item \textbf{Theoretical Concepts of Machine Learning:} estimation theory,
        VC-dimension, generalization bounds, convex optimization, Bayes theory.
\end{itemize}

The goal is not only to know individual methods, but to be able to \emph{choose}
the right approach for a given problem, understand the mathematical foundations well
enough to adapt standard algorithms, and recognize when a model is likely to fail.


\section{Course Specific Introduction}

A learning problem is called \emph{unsupervised} when the desired output is not given
for the training data.
In contrast to supervised learning, quality on unsupervised problems is measured as a
\emph{cumulative} quantity over all samples --- there is no per-sample ground truth.
Modern Generative AI (VAEs, GANs, diffusion models) is a prominent subtopic.

\paragraph{Objectives} Unsupervised methods are driven by different optimality criteria,
depending on the problem:
\begin{itemize}
  \item information content (maximizing information preserved after projection)
  \item orthogonality of extracted components
  \item statistical independence of components
  \item explained variance
  \item entropy of the representation
  \item likelihood: probability that the model produces the observed data
  \item distances within and between clusters
\end{itemize}

\paragraph{Use cases} Unsupervised methods are used to:
explore data, find latent structure, visualize high-dimensional data in low dimensions,
and compress representations.
The unifying goal is to \emph{understand and explore the data and generate new knowledge}.

\paragraph{Key concepts in this course}
Maximum likelihood (ML), maximum a posteriori (MAP), maximum entropy, expectation
maximization (EM), maximal variance, statistical independence, non-Gaussianity,
sub- and super-Gaussian distributions, sparse codes and population codes.

\paragraph{Model selection and overfitting}
The goal of learning is to select, from a model class, the model with the highest
\emph{generalization performance} --- the model that best explains \emph{future} data.
Model selection, training, and learning are all the same thing.
\emph{Overfitting} occurs when the model is fitted to idiosyncratic training
characteristics (noise, outliers, labeling errors) rather than to the true underlying
structure; the model has high training performance but poor generalization.

\begin{intbox}
  Underfitting (model too simple) and overfitting (model too complex) are the two
  failure modes. The goal is the best bias-variance trade-off: a model complex enough
  to capture the true structure but not so complex that it memorizes noise.
\end{intbox}

\paragraph{Parametric vs.\ nonparametric models}
An important early decision is the \emph{model class}:

\begin{itemize}
  \item \textbf{Parametric models:} each parameter vector $\wvec$ represents one model.
        Examples: neural networks (synaptic weights), SVMs (weight vector $\wvec$), factor
        analysis. Learning corresponds to a path through parameter space.
        Disadvantages: different parameterizations can represent the same function; model
        complexity is implicitly controlled through the parameters.

  \item \textbf{Nonparametric models:} the model is locally constant or a superimposition
        of local pieces. Examples: $k$-nearest-neighbor (where $k$ is a hyperparameter fixed
        before training), kernel density estimation, decision trees.
        The constant rules/kernels must be chosen a priori; hyperparameters ($k$, kernel
        width, splitting rules) are not learned from data.
\end{itemize}

\paragraph{Unsupervised methods covered in this course}
Principal component analysis (PCA), independent component analysis (ICA), factor analysis,
projection pursuit, $k$-means clustering, hierarchical clustering, mixture models
(Gaussian, Poisson), self-organizing maps (SOMs), kernel density estimation, hidden Markov
models (HMMs, factorial HMMs, IO-HMMs), Markov networks, restricted Boltzmann machines,
auto-associators.

Methods fall into two broad families:
\begin{itemize}
  \item \textbf{Projection methods:} produce a new representation of objects by
        down-projecting feature vectors into a lower-dimensional space while preserving
        neighborhoods or other desired properties.
        Examples: PCA (orthogonal components of maximal variance), ICA (statistically
        independent components), factor analysis (PCA with noise).
  \item \textbf{Generative models:} build an explicit probabilistic model of the observed
        data; the fitted model should match the observed data density.
\end{itemize}


\section{Generative vs.\ Descriptive Models}
\label{sec:gen-vs-desc}

\begin{defbox}[title={Descriptive (Recoding) Model}]
A \textbf{descriptive model} provides an alternative representation or description of the
data. It maps observations to a new space with desired properties (e.g.\ orthogonality,
independence). Examples: PCA, ICA.
\end{defbox}

\begin{defbox}[title={Generative Model}]
A \textbf{generative model} explicitly models the data generation process.
It should produce the same distribution as the observed data: random source components
$\uvec$ are drawn from a prior $p_{\uvec}$, passed through a model
$g(\uvec;\wvec)$, and the resulting distribution $p(\xvec;\wvec)$ should match the
empirical distribution $p(\xvec)$.
Examples: factor analysis, latent variable models, Boltzmann machines, HMMs.
\end{defbox}

Generative models allow the use of prior knowledge about the world and enable the
prediction of new states of the data-generating process (e.g.\ predicting gene expression
or brain activity). The natural training objective is maximum likelihood or MAP.

Descriptive models aim to find a richer or more interpretable representation rather than
to model the generative process.
Their objectives are properties of the representation: maximal variance (PCA),
statistical independence (ICA), maximal non-Gaussianity (projection pursuit).

\begin{notebox}[title={Exam trap}]
Parametric/nonparametric is \emph{not} the same as generative/descriptive.
Factor analysis is parametric \emph{and} generative.
$k$-NN is nonparametric \emph{and} descriptive (no explicit density).
\end{notebox}


% ==========================================================
\chapter{Basic Terms and Concepts}
% ==========================================================

\section{Unsupervised Learning in Bioinformatics}

The primary driving application for many unsupervised methods is gene expression data.
A microarray experiment produces a matrix of expression values: rows are genes (features),
columns are samples (conditions, patients, time points).
The data are high-dimensional (thousands of genes), noisy, and unlabeled.

Typical tasks include:
\begin{itemize}
  \item finding groups of co-expressed genes (clustering, biclustering)
  \item identifying tissue types or disease subtypes (mixture models, ICA)
  \item visualizing the data in 2--3 dimensions (PCA, $t$-SNE, Isomap)
  \item detecting regulatory modules that activate in a subset of samples (FABIA)
\end{itemize}

A classical result is Spellman's cell-cycle data: PCA applied to yeast expression data
over the cell cycle reveals the dominant temporal oscillation in the first two principal
components, confirming the cyclic structure.
Hierarchical clustering of microarray data yields dendrograms that group both genes and
samples; the correlation coefficient bar on the dendrogram indicates relatedness of branches.


\section{Unsupervised Learning Categories}

All unsupervised methods fall into one of two frameworks:

\begin{defbox}[title={Two Frameworks of Unsupervised Learning}]
\begin{itemize}
  \item \textbf{Generative framework:} density estimation, hidden Markov models, mixture
        models. Objective: \emph{maximum likelihood} or \emph{maximum a posteriori}.
  \item \textbf{Recoding (descriptive) framework:} projection methods, PCA, ICA.
        Objective: \emph{maximal variance}, orthogonality, independence, \emph{maximum entropy}.
\end{itemize}
\end{defbox}

\subsection{Generative Framework}

In the generative framework, the world is modeled as a stochastic process:
latent (hidden) sources $\uvec$ are drawn from a source distribution $p_{\uvec}$,
passed through a world model $g(\uvec;\wvec)$, and produce observations $\xvec$.
The learned model distribution $p(\xvec;\wvec)$ must match the observed distribution $p(\xvec)$.

\begin{thmbox}[title={Generative Model Objective}]
Find parameters $\wvec$ such that the model distribution
\[
  p(\xvec;\wvec) = \int_U p_{\uvec}(\uvec)\,\delta\!\bigl(\xvec = g(\uvec;\wvec)\bigr)\dd\uvec
\]
is as close as possible to the empirical distribution $p(\xvec)$.
The loss function measures the distance between these two distributions; typically this
is minimized via \emph{maximum likelihood}.
\end{thmbox}

Key examples of generative models: factor analysis, latent variable models,
Boltzmann machines, hidden Markov models.

\subsection{Recoding Framework}

In the recoding (descriptive) framework, a model $g(\xvec;\wvec)$ maps observations
$\xvec \sim p(\xvec)$ to projections $\uvec$.
The resulting distribution $p(\uvec;\wvec)$ should match a desired \emph{target}
distribution $p_t(\uvec)$:

\begin{itemize}
  \item \textbf{PCA:} projects to a low-dimensional space with maximal information
        conservation (maximal variance). Target: components are orthogonal and uncorrelated.

  \item \textbf{ICA:} projects into a space with statistically \emph{independent}
        components (factorial code). Often characteristics of a factorial distribution are
        optimized: maximal entropy given fixed variance, or cumulants (kurtosis).
        Alternatively, the projection distribution is matched to a product of
        super-Gaussians.

  \item \textbf{Projection Pursuit:} projects such that components are maximally
        non-Gaussian (since Gaussian components arise from random projections of any
        distribution, by the CLT; non-Gaussian components are informative).
\end{itemize}

\subsection{Recoding and Generative Framework Unified}

Both frameworks can be seen as two sides of the same coin.
In the generative direction: sources $\uvec \xrightarrow{g(\uvec;\wvec)} \xvec$
(sources produce observations).
In the recoding direction: observations $\xvec \xrightarrow{g^{-1}(\xvec;\wvec)} \uvec$
(observations are encoded into sources).
When the model $g$ is invertible, these two views are mathematically equivalent and
yield the same parameter estimates.
This unification is exploited in ICA: the INFOMAX and likelihood approaches give
identical learning rules for invertible linear models.


\section{Quality of Parameter Estimation}
\label{sec:quality-estimation}

Generative models must estimate true (but unknown) model parameters.
Let the training data be $\{\xvec\} = \{\xvec^1,\ldots,\xvec^l\}$, collected into the
matrix $\Xmat = (\xvec^1,\ldots,\xvec^l)\T \in \R^{l \times d}$.
The true (unknown) parameter vector is $\wvec$; an estimator maps the data to $\what$.

\begin{defbox}[title={Estimator Quality Measures}]
\begin{align*}
  \text{Unbiased:} \quad & \E_{\Xmat}[\what] = \wvec \\[4pt]
  \text{Bias:} \quad & \bias(\what) = \E_{\Xmat}[\what] - \wvec \\[4pt]
  \text{Variance:} \quad & \Var(\what) = \E_{\Xmat}\!\left[(\what - \E_{\Xmat}[\what])\T(\what - \E_{\Xmat}[\what])\right] \\[4pt]
  \text{MSE:} \quad & \mse(\what) = \E_{\Xmat}\!\left[(\what - \wvec)\T(\what - \wvec)\right]
\end{align*}
\end{defbox}

The fundamental bias-variance decomposition of the MSE is:

\begin{thmbox}[title={Bias-Variance Decomposition}]
\[
  \mse(\what) = \Var(\what) + \bias^2(\what)
\]
\emph{Derivation sketch:} Write $\what - \wvec = (\what - \E[\what]) + (\E[\what] - \wvec)$,
expand the squared norm, and note that the cross-term vanishes because
$\E_{\Xmat}[(\what - \E[\what])(\E[\what] - \wvec)\T] = 0$
(since $\E[\what] - \wvec$ is a constant w.r.t.\ the expectation).
\end{thmbox}

The objective is to minimize the MSE.
Note: the MSE in the parameter estimation sense (expected squared distance to true
parameter) is different from the supervised regression loss (expected squared distance
to true output).

\begin{intbox}
  Unbiasedness and consistency are different properties.
  An \emph{unbiased} estimator has $\E[\what] = \wvec$ for \emph{any} dataset size.
  A \emph{consistent} estimator satisfies $\what \to \wvec$ as $l \to \infty$ ---
  its variance shrinks with more data, but it may be biased for finite $l$.
\end{intbox}


\section{Maximum Likelihood Estimator}
\label{sec:mle}

\begin{defbox}[title={Likelihood}]
The \textbf{likelihood} of the dataset $\{\xvec\} = \{\xvec^1,\ldots,\xvec^l\}$ under
model parameters $\wvec$ is
\[
  \lhood(\{\xvec\};\wvec) \;=\; p(\{\xvec\};\wvec),
\]
the probability that the model $p(\xvec;\wvec)$ produces the observed data.
For \textbf{iid} (independent, identically distributed) data this factorizes:
\[
  \lhood(\{\xvec\};\wvec) \;=\; \prod_{i=1}^{l} p(\xvec^i;\wvec).
\]
Working in log-space converts the product to a sum (negative log-likelihood, NLL):
\[
  -\ln \lhood(\{\xvec\};\wvec) \;=\; -\sum_{i=1}^{l} \ln p(\xvec^i;\wvec).
\]
\end{defbox}

Maximum likelihood estimation (MLE) is the dominant objective in unsupervised learning.
It finds the parameters $\wvec$ that make the observed data most probable under the model:
\[
  \what_{\mathrm{ML}} = \argmax_{\wvec} \lhood(\{\xvec\};\wvec) = \argmin_{\wvec} \bigl[-\ln\lhood(\{\xvec\};\wvec)\bigr].
\]

\begin{thmbox}[title={Properties of the MLE}]
The maximum likelihood estimator is:
\begin{itemize}
  \item \textbf{Invariant under reparameterization:} if $\hat{\wvec}$ is the MLE of
        $\wvec$, then $f(\hat{\wvec})$ is the MLE of $f(\wvec)$ for any function $f$.
  \item \textbf{Asymptotically unbiased and efficient:} for large $l$, the MLE achieves
        the Cram\'{e}r-Rao lower bound, i.e.\ it has the smallest possible variance among
        all unbiased estimators. It is asymptotically optimal.
  \item \textbf{Asymptotically consistent:} $\what \to \wvec$ as $l \to \infty$.
\end{itemize}
\end{thmbox}

\begin{notebox}[title={MLE caveats}]
MLE can overfit for small datasets (it uses no prior information).
For finite $l$, MLE may be biased: the classic example is the MLE of the variance of a
Gaussian, $\hat{\sigma}^2 = \frac{1}{l}\sum(x_i - \bar{x})^2$, which underestimates
the true variance (the unbiased estimator uses $l-1$).
\end{notebox}


\section{Expectation Maximization}
\label{sec:em}

Many generative models include \emph{latent (hidden) variables} $\uvec$ that are not
directly observed.
The marginal likelihood then requires integrating over all latent configurations:
\[
  p(\xvec;\wvec) = \int_U p_{\uvec}(\uvec)\,\delta\!\bigl(\xvec = g(\uvec;\wvec)\bigr)\dd\uvec.
\]
This integral is in general intractable due to:
\begin{itemize}
  \item hidden states (discrete or continuous latent variables),
  \item many-to-one output mappings ($g$ is not injective),
  \item non-linearities in $g$.
\end{itemize}

The \textbf{Expectation Maximization (EM)} algorithm circumvents this by working with
the \emph{joint} distribution $p(\xvec, \uvec;\wvec)$, which is easier to handle, and
by constructing a tractable lower bound on the log-likelihood.

\paragraph{Derivation of the lower bound}
Introduce an arbitrary distribution $Q(\uvec \given \{\xvec\})$ over latent variables.
Apply Jensen's inequality to the concave logarithm:
\begin{align*}
  \ln \lhood(\{\xvec\};\wvec)
  &= \ln \int_U p(\{\xvec\},\uvec;\wvec)\,\dd\uvec \\
  &= \ln \int_U Q(\uvec \given \{\xvec\})\,
     \frac{p(\{\xvec\},\uvec;\wvec)}{Q(\uvec \given \{\xvec\})}\,\dd\uvec \\
  &\geq \int_U Q(\uvec \given \{\xvec\})\,
     \ln\frac{p(\{\xvec\},\uvec;\wvec)}{Q(\uvec \given \{\xvec\})}\,\dd\uvec
  \;\eqdef\; \Fbound(Q,\wvec).
\end{align*}
$\Fbound(Q,\wvec)$ is the \textbf{free energy} (ELBO: evidence lower bound).

\begin{thmbox}[title={EM Algorithm}]
EM alternates between two steps until convergence:
\begin{align*}
  \textbf{E-step:} \quad & Q_{k+1} = \argmax_{Q}\; \Fbound(Q,\wvec_k)
    \quad \Rightarrow \quad Q_{k+1} = p(\uvec \given \{\xvec\};\wvec_k) \\[4pt]
  \textbf{M-step:} \quad & \wvec_{k+1} = \argmax_{\wvec}\; \Fbound(Q_{k+1},\wvec)
\end{align*}
\end{thmbox}

\paragraph{Monotone convergence}
Both steps can only increase (or leave unchanged) the lower bound $\Fbound$.
The E-step tightens the bound: at the beginning of the M-step,
$\Fbound(Q_{k+1},\wvec_k) = \ln\lhood(\{\xvec\};\wvec_k)$ (the bound is tight).
The M-step then increases $\Fbound$, which implies:
\[
  \ln\lhood(\{\xvec\};\wvec_k) \;\leq\; \Fbound(Q_{k+1},\wvec_{k+1}) \;\leq\; \ln\lhood(\{\xvec\};\wvec_{k+1}).
\]
Therefore the log-likelihood is guaranteed to be non-decreasing at every EM iteration.

\begin{intbox}
  The E-step does not change model parameters --- it only updates the posterior
  distribution over latent variables. The M-step maximizes the expected complete-data
  log-likelihood, which for exponential family models has a closed-form solution.
  EM converges to a local maximum (not necessarily global).
\end{intbox}

EM is applied throughout this course: Baum-Welch algorithm for HMMs, fitting Gaussian
mixtures, factor analysis, and (soft) ICA variants.


\section{Maximum Entropy}
\label{sec:maxent}

\begin{defbox}[title={Shannon Entropy (discrete case)}]
For a discrete distribution $(p_1, p_2, \ldots)$:
\[
  H = -\sum_{k \geq 1} p_k \log p_k,
  \qquad \text{with the convention } 0 \cdot \log 0 = 0.
\]
\end{defbox}

The \textbf{principle of maximum entropy} (Jaynes, 1957) states: among all distributions
consistent with the available constraints (e.g.\ known moments), choose the one with
\emph{maximum entropy}.
This distribution:
\begin{itemize}
  \item makes the minimal assumptions beyond the stated constraints,
  \item corresponds to the most likely empirical distribution that could give rise to the
        observed features,
  \item connects statistical mechanics (thermodynamic equilibrium) and information theory.
\end{itemize}

\begin{notebox}[title={Existence caveat}]
Not every class of distributions contains a maximum entropy member.
For example, the class of all distributions with a given mean has arbitrarily large
entropy (entropy can be pushed to $\infty$).
The class of distributions with mean zero, second moment one, and third moment one has
entropy bounded from above but the supremum is not attained.
\end{notebox}

\paragraph{Maximum entropy with moment constraints}
Given moment constraints $\sum_{i=1}^n p(x_i)\,f_k(x_i) = F_k$ for $k=1,\ldots,m$
and normalization $\sum_{i=1}^n p(x_i) = 1$, the maximum entropy solution is the
\textbf{Gibbs distribution}:

\begin{thmbox}[title={Gibbs Distribution (MaxEnt Solution)}]
\[
  p(x_i) = \frac{1}{Z(\lambda_1,\ldots,\lambda_m)}
  \exp\!\bigl(\lambda_1 f_1(x_i) + \cdots + \lambda_m f_m(x_i)\bigr),
\]
with \textbf{partition function}
\[
  Z(\lambda_1,\ldots,\lambda_m) = \sum_{i=1}^n \exp\!\bigl(\lambda_1 f_1(x_i) + \cdots + \lambda_m f_m(x_i)\bigr).
\]
The Lagrange multipliers $\lambda_k$ are determined by the constraint equations:
\[
  F_k = \pd{}{\lambda_k} \log Z(\lambda_1,\ldots,\lambda_m), \qquad k = 1,\ldots,m.
\]
\end{thmbox}

\begin{exbox}[title={Standard MaxEnt examples}]
\begin{itemize}
  \item \textbf{Normal distribution:} maximum entropy distribution given fixed mean and
        variance. Constraints: $\E[X] = \mu$ and $\E[X^2] = \mu^2 + \sigma^2$.
  \item \textbf{Uniform distribution:} maximum entropy on a bounded interval $[a,b]$.
        Only constraint: support.
  \item \textbf{Exponential distribution:} maximum entropy on $[0,\infty)$ given fixed mean.
        Constraint: $\E[X] = 1/\lambda$.
\end{itemize}
\end{exbox}

\begin{intbox}
  The MaxEnt principle explains why Gaussian distributions appear so frequently in
  nature: they are the maximum entropy distributions for processes with finite variance.
  Any physical system evolving under energy constraints (finite second moment) will
  converge to a Gaussian-like distribution.
  In ICA, maximizing entropy of a nonlinear transform of the data is equivalent to
  maximizing likelihood of an independent-component model (INFOMAX).
\end{intbox}


% ==========================================================
\chapter{Principal Component Analysis}
% ==========================================================

\section{The Method}

Let the data matrix be $\bm{X} \in \R^{n \times m}$, where $n$ is the number of samples
and $m$ is the number of features (the rows $\bm{x}_i \in \R^m$ are centered so that
$\sum_i \bm{x}_i = \bm{0}$).

The \textbf{sample covariance matrix} is
\[
  \bm{C} = \frac{1}{n}\bm{X}\T\bm{X} \in \R^{m \times m}.
\]
Its eigendecomposition $\bm{C} = \frac{1}{n}\bm{U}\bm{D}_m\bm{U}\T$ yields an orthogonal
matrix $\bm{U}$ whose columns $\bm{u}_k$ are the principal directions (\emph{loadings}),
and a diagonal matrix $\bm{D}_m = \mathrm{diag}(\lambda_1,\ldots,\lambda_m)$ of
eigenvalues sorted in decreasing order.

The \textbf{SVD} of $\bm{X}$ gives the same result: $\bm{X} = \bm{V}\bm{D}\bm{U}\T$
with $\bm{V} \in \R^{n \times n}$, $\bm{D} \in \R^{n \times m}$ diagonal, and
$\bm{U} \in \R^{m \times m}$ orthogonal.
The eigenvalues of $\bm{C}$ are $\lambda_k = d_k^2/n$ where $d_k$ are the singular values.

\begin{defbox}[title={PCA Projection}]
The \textbf{principal component scores} (projections) are collected in
\[
  \bm{Y} = \bm{X}\bm{U} = \bm{V}\bm{D} \in \R^{n \times m}.
\]
The $k$-th column $\bm{y}_k = \bm{X}\bm{u}_k$ is the projection of all data points
onto the $k$-th principal direction.
The original data can be reconstructed as the outer product sum
\[
  \bm{X} = \sum_{k=1}^{m} \bm{y}_k \bm{u}_k\T.
\]
\end{defbox}

\paragraph{Online learning: Oja's rule}
A single-component PCA can be learned online (one sample at a time) using
\textbf{Oja's rule}:
\[
  \bm{u}^{\mathrm{new}} = \bm{u} + \eta\bigl(t\,\bm{x} - t^2\,\bm{u}\bigr),
  \qquad t = \bm{u}\T\bm{x}.
\]
This is a fixed-point iteration analogous to the power method.


\section{Variance Maximization}

PCA can be motivated as the projection that maximizes the variance of the projected data.

The first principal direction solves
\[
  \bm{u}_1 = \argmax_{\|\bm{u}\|=1} \sum_{i=1}^{n}(\bm{u}\T\bm{x}_i)^2
  = \argmax_{\|\bm{u}\|=1} \bm{u}\T\bm{C}\bm{u}.
\]

\begin{thmbox}[title={Variance Maximization Solution}]
Using the eigendecomposition $\bm{C} = \sum_k \lambda_k \bm{u}_k\bm{u}_k\T$ and writing
$\bm{u} = \sum_k a_k \bm{u}_k$, the objective becomes
$\bm{u}\T\bm{C}\bm{u} = \sum_k a_k^2 \lambda_k$ subject to $\sum_k a_k^2 = 1$.
This is maximized by $a_1 = 1$, all others zero:
\[
  \bm{u}_1 = \text{eigenvector of } \bm{C} \text{ with largest eigenvalue } \lambda_1.
\]
The $k$-th principal direction is the direction of maximal variance \emph{orthogonal} to
all previous $k-1$ directions.
The variance explained by the $k$-th component equals $\lambda_k$.
\end{thmbox}


\section{Uniqueness}

PCA is \textbf{unique up to sign flips} of the principal directions, provided all
eigenvalues are distinct.
If two or more eigenvalues are equal (degenerate case), the corresponding eigenvectors
span a subspace but their individual directions are arbitrary; any orthonormal basis
of that subspace is equally valid.

\begin{notebox}[title={Zero eigenvalue}]
At most one eigenvalue can be zero (a zero eigenvalue means the data lies in a
lower-dimensional affine subspace; in that case the corresponding feature can be dropped
before running PCA).
\end{notebox}


\section{Properties of PCA}
\label{sec:pca-properties}

Let $\bm{y}^{(i)} = \bm{U}\T\bm{x}_i$ be the projection of sample $i$.

\begin{thmbox}[title={Key Properties of PCA Projections}]
\begin{itemize}
  \item \textbf{Zero mean:} $\sum_i y_k^{(i)} = 0$ for all $k$ (because the data are centered).
  \item \textbf{Uncorrelated:} $\mathrm{Cov}(y_k, y_l) = 0$ for $k \neq l$
        (follows from orthogonality of $\bm{U}$).
  \item \textbf{Variance:} $\mathrm{Var}(y_k) = \lambda_k$ (the $k$-th eigenvalue).
  \item \textbf{Optimal reconstruction:} truncating to the first $l$ components gives
        the best rank-$l$ approximation in terms of MSE:
        \[
          \mathrm{MSE}_l = \sum_{k=l+1}^{m} \lambda_k.
        \]
  \item \textbf{Fraction of variance} explained by the first $l$ components:
        $\frac{\sum_{k=1}^l \lambda_k}{\sum_{k=1}^m \lambda_k}$.
\end{itemize}
\end{thmbox}


\section{Examples}

\subsection{Iris Data Set}

The Iris dataset has $n = 150$ samples (50 per species: \textit{setosa}, \textit{versicolor},
\textit{virginica}) and $m = 4$ features (sepal length, sepal width, petal length, petal width).

PC1 accounts for approximately \textbf{92\% of variance} and is the only component that
separates the three species.
Pairwise scatter plots of all four PCs show that only plots involving PC1 reveal a
clear three-group structure; all other pairs (PC2 vs.\ PC3, etc.)\ show no meaningful
separation.
This tells us the discriminative signal is essentially one-dimensional.

\subsection{Multiple Tissue Data Set}

This microarray dataset contains expression values for 5{,}565 genes across 102 samples
from four human tissue types: \textbf{breast} (Br), \textbf{prostate} (Pr),
\textbf{lung} (Lu), and \textbf{colon} (Co).

Applied to all genes simultaneously:
\begin{itemize}
  \item PC1 (5\% variance) separates prostate from the rest.
  \item PC2 (3\%) separates colon and also breast samples.
  \item PC3 (3\%) partially separates lung.
\end{itemize}

\paragraph{Variance filtering}
Selecting only the top-variance genes before PCA is justified for microarray data
(most genes are essentially constant noise).
Three filtered subsets were compared:
\begin{itemize}
  \item \textbf{101 genes} (highest variance, XMultiF1): PC1 (36\%) separates prostate;
        PC2 (14\%) separates colon; PC3 (8\%) separates breast and lung.
        Much cleaner separation than the full set.
  \item \textbf{13 genes} (XMultiF2, 57\% PC1): similar tissue-level separation.
  \item \textbf{5 genes} (XMultiF3, 77\% PC1): still separates prostate clearly but
        other tissues become hard to distinguish.
        Inspection reveals 4 of the 5 selected genes are highly correlated (ACPP, KLK2,
        MSMB, TRGC2 --- all prostate-related; KRT5 is not).
\end{itemize}

\paragraph{Gene clustering}
Applying hierarchical clustering with correlation as distance metric, then selecting the
gene with maximal variance per cluster (\textbf{92 uncorrelated genes of maximal variance})
gives very similar results to variance-based filtering.
The approach ensures genes are not redundant (correlated genes contribute the same
information to PCA) while still selecting high-variance representatives.


\section{Kernel Principal Component Analysis}

\subsection{The Method}

\textbf{Kernel PCA} (KPCA) extends PCA to nonlinear projections by performing the
linear operations of PCA inside a \emph{reproducing kernel Hilbert space} (RKHS), to
which the data are first mapped non-linearly:
\[
  \bm{x} \;\mapsto\; \Phi(\bm{x}).
\]

Assuming the mapped data are centered ($\sum_i \Phi(\bm{x}_i) = \bm{0}$), the
covariance in feature space is
$\bm{C} = \frac{1}{n}\sum_i \Phi(\bm{x}_i)\Phi(\bm{x}_i)\T$,
and the Gram (kernel) matrix has entries $K_{ij} = \Phi(\bm{x}_j)\T\Phi(\bm{x}_i) = k(\bm{x}_i,\bm{x}_j)$.

The eigenvalue problem $\bm{C}\bm{w} = \lambda\bm{w}$ cannot be solved in feature space
directly (we only have $K_{ij}$), but the solution must lie in the span of the mapped
data:
\[
  \bm{w} = \sum_{i=1}^n \alpha_i\,\Phi(\bm{x}_i).
\]
Substituting and using $K_{ij}$ yields the \textbf{kernel eigenvalue problem}:
\[
  n\lambda\,\bm{K}\bm{\alpha} = \bm{K}^2\bm{\alpha}
  \;\;\Rightarrow\;\;
  n\lambda\,\bm{\alpha} = \bm{K}\bm{\alpha}.
\]
Normalization $\|\bm{w}\|=1$ requires $\|\bm{\alpha}\| = 1/\sqrt{\tilde\lambda}$ where
$\tilde\lambda = n\lambda$, so
\[
  \bm{\alpha} = \frac{\tilde{\bm{\alpha}}}{\sqrt{\tilde\lambda}}.
\]

The projection of a new point $\bm{x}$ onto a kernel principal direction is then
\[
  \bm{w}\T\Phi(\bm{x}) = \sum_{i=1}^n \alpha_i\,k(\bm{x}_i,\bm{x}).
\]

\begin{thmbox}[title={Kernel PCA Algorithm}]
\begin{enumerate}
  \item \textbf{Center} the Gram matrix:
        $\tilde{\bm{K}} = \bm{K} - \frac{1}{n}\bm{K}\bm{1}\bm{1}\T - \frac{1}{n}\bm{1}\bm{1}\T\bm{K} + \frac{1}{n^2}(\bm{1}\T\bm{K}\bm{1})\bm{1}\bm{1}\T$.
  \item \textbf{Eigenvalues}: compute eigenvectors $\bm{\alpha}$ and eigenvalues
        $\tilde\lambda$ of $\tilde{\bm{K}}$.
  \item \textbf{Normalize}: $\alpha_i^{\mathrm{new}} = \alpha_i/(\|\bm{\alpha}\|\sqrt{n\lambda})$.
  \item \textbf{Project} a new vector $\bm{x}$: center its kernel vector and compute
        $\bm{w}\T\Phi(\bm{x}) = \sum_i \alpha_i\,k(\bm{x}_i,\bm{x})$.
\end{enumerate}
\end{thmbox}

\begin{intbox}
  Kernel PCA can detect nonlinear structure that linear PCA misses.
  For example, concentric rings are linearly inseparable but perfectly separable in
  the RKHS of an RBF kernel.
  The price is that the representation depends on all training points (no compact weight
  vector $\bm{u}$), and choosing the right kernel function requires domain knowledge.
\end{intbox}


% ==========================================================
\chapter{Independent Component Analysis}
% ==========================================================

Independent component analysis (ICA) seeks a representation of the data as a linear
mixture of \textbf{statistically independent} non-Gaussian sources.

\begin{defbox}[title={ICA Model}]
\textbf{Generative (mixing) model:}
\[
  \bm{x} = \bm{U}\,\bm{y}, \qquad p(\bm{y}) = \prod_{j=1}^l p(y_j).
\]
$\bm{U} \in \R^{m \times l}$ is the mixing matrix; the sources $y_j$ are statistically
independent.

\textbf{Descriptive (demixing) model:}
\[
  \bm{y} = \bm{W}\,\bm{x}, \qquad \bm{W} = \bm{U}^{-1}
  \quad (\text{for } l = m).
\]
Matrix form: $\bm{X} = \bm{Y}\bm{U}\T$, with $\bm{Y}\T\bm{Y} = \bm{D}_m$ (decorrelated,
and in fact statistically independent).
Unlike PCA, $\bm{U}$ is \textbf{not required to be orthogonal}.
\end{defbox}

The canonical motivation is \textbf{blind source separation}: two speakers talk
simultaneously; multiple microphones record linear mixtures of the two voices.
ICA recovers the original signals from the mixtures without knowing the mixing matrix.


\section{Identifiability and Uniqueness}
\label{sec:ica-uniqueness}

The ICA solution is generally \textbf{not fully unique}: $\bm{x} = \bm{U}\bm{P}^{-1}\bm{P}\bm{y}$
gives another valid decomposition for any invertible $\bm{P}$.
However, if $\bm{P}$ cannot mix the statistically independent components of $\bm{y}$,
then $\bm{P}$ must be a product of a permutation matrix and a diagonal scaling matrix.

\begin{thmbox}[title={Darmois' Theorem (1953) --- ICA Identifiability}]
Define $x_1 = \sum_j a_j y_j$ and $x_2 = \sum_j b_j y_j$ where $y_j$ are independent
random variables.
If $x_1$ and $x_2$ are independent, then all $y_j$ for which $a_j b_j \neq 0$ are
\emph{Gaussian}.

\medskip
Consequence: for \textbf{non-Gaussian sources}, the ICA solution is unique up to
permutation and scaling.
The Gaussian case is the only exception.
\end{thmbox}

\paragraph{ICA assumptions}
\begin{itemize}
  \item Sources are \textbf{non-Gaussian} (at most one source can be Gaussian).
  \item $l \leq m$ (at least as many observations as sources).
  \item $\bm{U}$ has full rank $l$.
\end{itemize}
Source densities $p(y_i)$ are typically approximated by super-Gaussian or unimodal
distributions (generative framework).

\paragraph{Two criteria for independence}
\begin{enumerate}
  \item \textbf{Mutual information} between components.
  \item \textbf{Non-Gaussianity} of components.
\end{enumerate}


\section{Measuring Independence}

\subsection{Mutual Information}

The \textbf{mutual information} of a vector-valued random variable measures how far the
joint distribution is from a factorial (product) distribution:
\[
  I(y_1,\ldots,y_l) = \sum_{j=1}^l H(y_j) - H(\bm{y}),
\]
where $H(\bm{a}) = -\int p(\bm{a})\ln p(\bm{a})\dd\bm{a}$ is the (differential) entropy.
$I \geq 0$, with equality if and only if all components are statistically independent.

For the linear demixing $\bm{y} = \bm{W}\bm{x}$, using the change-of-variables formula
$p(\bm{y}) = p(\bm{x})/|\bm{W}|$:
\[
  I(y_1,\ldots,y_m) = \sum_{j=1}^m H(y_j) - H(\bm{x}) - \ln|\bm{W}|.
\]
Since $H(\bm{x})$ is constant in $\bm{W}$, minimizing mutual information amounts to
maximizing $\ln|\bm{W}| - \sum_j H(y_j)$, which is the ICA objective.

\subsection{Non-Gaussianity}

\paragraph{Negentropy}
The \textbf{negentropy} measures how non-Gaussian a distribution is:
\[
  J(\bm{y}) = H(\bm{y}_{\mathrm{gauss}}) - H(\bm{y}),
\]
where $\bm{y}_{\mathrm{gauss}}$ is a Gaussian with the same covariance as $\bm{y}$.
$J \geq 0$ (Gaussian has maximum entropy for fixed variance); maximizing negentropy
minimizes mutual information.
Direct estimation of negentropy is difficult in practice.

\paragraph{Kurtosis}
Non-Gaussianity can be quantified via cumulants, in particular the
\textbf{kurtosis} (fourth cumulant):
\[
  \kappa_4 = \E(x^4) - 3\bigl(\E(x^2)\bigr)^2.
\]
For a Gaussian, $\kappa_3 = \kappa_4 = 0$.
\begin{itemize}
  \item $\kappa_4 > 0$: \textbf{super-Gaussian} (heavier tails than Gaussian, e.g.\
        Laplace distribution).
  \item $\kappa_4 < 0$: \textbf{sub-Gaussian} (lighter tails, e.g.\ uniform distribution).
\end{itemize}
Key properties:
\begin{itemize}
  \item Additivity: $\kappa_4(x_1 + x_2) = \kappa_4(x_1) + \kappa_4(x_2)$ for
        independent $x_1, x_2$.
  \item Scaling: $\kappa_4(\alpha x) = \alpha^4 \kappa_4(x)$.
  \item Mixing \emph{reduces} kurtosis toward zero (towards Gaussian by CLT).
        Therefore \textbf{maximizing kurtosis} recovers the unmixed sources.
\end{itemize}
Kurtosis is not robust: as a fourth-order moment, it is heavily influenced by outliers.

\paragraph{Sparseness}
Maximizing kurtosis is equivalent to maximizing \textbf{sparseness}: a sparse variable
is almost always near zero but occasionally takes large values (relative to a Gaussian
with the same variance).
Sparseness does \emph{not} imply small variance.

\paragraph{Contrast functions}
Several contrast functions are used in practice to measure non-Gaussianity:
\begin{itemize}
  \item Kurtosis: $\kappa_4(y)$.
  \item Mixed cumulants: $\frac{1}{12}\kappa_3^2(y) + \frac{1}{48}\kappa_4^2(y)$
        (normalized to zero mean, unit variance).
  \item $\left|\E_y[G(y)] - \E_\nu[G(\nu)]\right|^p$ where $\nu$ is a standard Gaussian
        and $G$ is a nonlinearity; robust choices include
        $G(x) = \log\cosh(ax)$ and $G(x) = \exp(-ax^2/2)$ with $a \geq 1$.
\end{itemize}


\section{Whitening and Rotation Algorithms}

ICA is not unique up to scaling, so the standard approach is to first remove the
scaling ambiguity by \textbf{whitening} (also called sphering), then find an
orthogonal \textbf{rotation} that maximizes non-Gaussianity.

\textbf{Whitened} data satisfies $\bm{Y}\T\bm{Y} = \bm{I}$.
Starting from $\bm{X} = \bm{Y}\bm{U}\T$ with $\bm{C} = \frac{1}{n}\bm{X}\T\bm{X}$:
\[
  \bm{I} = \bm{C}^{-1/2}\,\frac{1}{n}\bm{X}\T\bm{X}\,\bm{C}^{-1/2}
  = \frac{1}{n}\bm{C}^{-1/2}\bm{U}\,\bm{Y}\T\bm{Y}\,\bm{U}\T\bm{C}^{-1/2}.
\]
This means $\hat{\bm{U}} = \frac{1}{\sqrt{n}}\bm{C}^{-1/2}\bm{U}$ is orthogonal, and
the ICA solution decomposes as:
\[
  \bm{Y} = \frac{1}{\sqrt{n}}\,\bm{X}\,\bm{C}^{-1/2}\,\hat{\bm{U}}.
\]
The two-step procedure is:
\begin{enumerate}
  \item \textbf{Whiten:} compute $\tilde{\bm{X}} = \bm{X}\bm{C}^{-1/2}$.
  \item \textbf{Rotate:} find the orthogonal matrix $\hat{\bm{U}}$ that maximizes
        the contrast function (super-Gaussianity / sparseness).
\end{enumerate}


\section{INFOMAX Algorithm}

\textbf{INFOMAX} (Bell \& Sejnowski, 1995) minimizes mutual information by maximizing
the entropy of a nonlinear transform of $\bm{y}$:
\[
  \bm{y} = \bm{W}\bm{x}, \qquad
  H(\bm{g}(\bm{y})) \text{ is maximized},
  \quad \bm{g}(\bm{y}) = (g(y_1),\ldots,g(y_l))\T.
\]
At maximum entropy the mutual information $I(g(y_1),\ldots,g(y_l)) = 0$, so the
components are independent.
The nonlinearity $g$ is chosen to approximate the CDF of the source distribution;
a common choice is $g(y_i) = \tanh(y_i)$ (super-Gaussian sources).

The entropy expands as
\[
  H(\bm{g}(\bm{y})) \approx H(\bm{x}) + \frac{1}{n}\sum_{i=1}^n\sum_{j=1}^l
  |\ln g'(y_{ij})| + \ln|\bm{W}|.
\]

\begin{thmbox}[title={INFOMAX Update Rules (Natural Gradient)}]
\textbf{tanh} ($g(y_j) = \tanh(y_j)$):
\[
  \Delta\bm{W} \;\propto\; \bigl(\bm{I} - 2\,\bm{g}(\bm{y})\,\bm{y}\T\bigr)\bm{W}
\]
\textbf{sigmoid} ($g(y_j) = 1/(1+e^{-y_j})$):
\[
  \Delta\bm{W} \;\propto\; \bigl(\bm{I} + (1 - 2\,\bm{g}(\bm{y}))\,\bm{y}\T\bigr)\bm{W}
\]
INFOMAX is equivalent to a generative approach using \textbf{maximum likelihood}.
\end{thmbox}

\begin{intbox}
  The natural gradient (premultiplication by $\bm{W}\T\bm{W}$) converts the
  standard gradient $(\bm{W}\T)^{-1} \pm \ldots\bm{x}\T$ into the more elegant
  update involving $\bm{y} = \bm{W}\bm{x}$.
  It also makes the algorithm \emph{equivariant}: performance does not depend on the
  scaling or conditioning of the input.
\end{intbox}


\section{EASI Algorithm}

\textbf{EASI} (Equivariant Adaptive Separation via Independence, Cardoso \& Laheld, 1996)
uses the update rule
\[
  \Delta\bm{W} \;\propto\; \bigl(\bm{I} - \bm{y}\bm{y}\T - \bm{g}(\bm{y})\bm{y}\T
  + \bm{y}\,\bm{g}\T(\bm{y})\bigr)\bm{W},
\]
where $\bm{g}$ is the same class of contrast functions as in INFOMAX (tanh or sigmoid).
EASI is designed for online learning and is also equivariant by construction.


\section{FastICA Algorithm}

\textbf{FastICA} (Hyvärinen, 1999) is probably the most widely used ICA algorithm.
It is a \emph{fixed-point} iteration (analogous to Oja's rule for PCA) that maximizes
non-Gaussianity as measured by a contrast function $g$ with derivative $g'$:

\begin{thmbox}[title={FastICA Update}]
\[
  \bm{w}^{\mathrm{new}} = \E\bigl[\bm{x}\,g(\bm{w}\T\bm{x})\bigr]
  - \E\bigl[g'(\bm{w}\T\bm{x})\bigr]\,\bm{w}.
\]
The new $\bm{w}^{\mathrm{new}}$ is then normalized to unit length.
\end{thmbox}

Key properties:
\begin{itemize}
  \item Whitening and rotation algorithm.
  \item Generalizes kurtosis maximization to arbitrary contrast functions $g$.
  \item Extended to extract multiple independent components simultaneously
        (deflation or symmetric extraction).
\end{itemize}


\section{ICA Extensions}

Beyond the basic square, invertible model several extensions exist:
\begin{itemize}
  \item \textbf{Generative approach:} assumptions on the source densities $p(y_i)$,
        sub-Gaussian sources with specific priors.
  \item \textbf{Non-linear ICA:} mixing function $\bm{x} = g(\bm{y})$ is nonlinear;
        solutions are often not unique without additional constraints.
  \item \textbf{Overcomplete ICA:} more sources than observations ($l > m$); the system
        is underdetermined but useful for sparse representations of signals.
  \item \textbf{Undercomplete ICA:} fewer sources than observations ($l < m$); reduces
        to a lower-dimensional model.
\end{itemize}


\section{ICA vs.\ PCA}

\begin{center}
\renewcommand{\arraystretch}{1.25}
\begin{tabular}{ll}
  \toprule
  \textbf{ICA} & \textbf{PCA} \\
  \midrule
  Seeks \emph{causes} of the data & Seeks \emph{geometrical abstractions} \\
  Statistically independent components & Decorrelated (orthogonal) components \\
  Explains super-Gaussian structure & Explains all variance \\
  Scale invariant & Not scale invariant \\
  Unique up to scaling and permutation & Unique (up to sign) \\
  Assumes super-Gaussian sources & No distributional assumptions \\
  Components not ranked & Components ranked by eigenvalue \\
  \bottomrule
\end{tabular}
\end{center}

\begin{intbox}
  PCA gives the best linear dimensionality reduction in terms of MSE; ICA finds the
  interpretation that is statistically most meaningful for independent source models.
  For Gaussian data both methods give the same result (any rotation of decorrelated
  Gaussians is equally independent), which is why ICA requires non-Gaussianity.
\end{intbox}


\section{Artificial ICA Examples}

Starting with 1{,}000 points from two independent \textbf{uniform distributions}
(uniform = sub-Gaussian, $\kappa_4 < 0$), the three processing stages can be visualized:

\begin{enumerate}
  \item \textbf{Original sources:} square-shaped scatter (axes-aligned).
  \item \textbf{Mixing} ($\bm{X} = \bm{Y}\bm{U}\T$): the square is rotated and sheared
        into a parallelogram/diamond shape --- the components are now dependent.
  \item \textbf{Whitening} ($\bm{X}\bm{C}^{-1/2}$): the shape becomes circular
        (unit covariance); correlation is removed but statistical dependence remains.
  \item \textbf{Rotation} (ICA): the circle is aligned back to a square, recovering
        the original independent sources.
\end{enumerate}

Both INFOMAX and FastICA recover the original square (up to permutation and sign).
For super-Gaussian sources (Laplace distribution), the shape is a four-pointed star,
and the same pipeline recovers it correctly.


\section{Real World ICA Examples}

\subsection{Iris Data Set}

On the Iris data, ICs are \textbf{ordered by their impact on the observations}
(column norms of the mixing matrix $\bm{U}$), \emph{not} by variance.
The first independent component explains approximately 90\% of the variance in the
data and likely encodes the \emph{size of the blossom} (a single dominant biological
factor).
Unlike PCA, ICA does not guarantee components are ranked by variance.

\subsection{Multiple Tissue Data Set}

Applied to the multiple tissue microarray data (102 samples, 4 tissue types), the
first three ICs provide biologically meaningful separations:
\begin{itemize}
  \item \textbf{IC1} (29\% variance): separates \emph{secretory/reproductive} tissue
        (prostate and breast) from \emph{internal organ} tissue (colon and lung) ---
        a biologically interpretable axis.
  \item \textbf{IC2} (28\%): separates prostate + lung from breast + colon.
  \item \textbf{IC3} (27\%): separates prostate + colon from breast + lung.
  \item All combinations of the first three ICs yield clean separations except for
        partial overlap between breast and lung samples.
\end{itemize}
Genes most correlated to IC1 include SERPINA7, LAMB3, \textbf{AR}, CCNG2, KLF5, CCL20,
SLC39A14, ATP1B1, GSTP1, LAD1.
The \textbf{androgen receptor} (AR) is the 3rd most related gene to IC1 --- consistent
with IC1 separating reproductive/secretory tissues, since androgen signaling drives
prostate gland differentiation and also plays a role in normal breast physiology.


\section{Kurtosis Maximization Results in Independent Components}

The theoretical justification for kurtosis-based ICA rests on a key inequality:
the kurtosis of a linear mixture of independent variables is always less than or equal
to the \emph{maximum} kurtosis of the individual sources.
Formally, for $x = \bm{a}\T\bm{y}$ with $\|\bm{a}\| = 1$ and $y_j$ independent:
\[
  \kappa_4(x) = \sum_j a_j^4\,\kappa_4(y_j) \;\leq\; \max_j |\kappa_4(y_j)|.
\]
Equality holds when one $|a_j| = 1$ and all others are zero --- i.e.\ when the
projection aligns exactly with one source.
Therefore \textbf{maximizing $|\kappa_4|$ over all unit-norm projections recovers a
single source component}.
Iterating (deflating each recovered component) separates all sources.

\begin{notebox}[title={Limitations of kurtosis}]
Because kurtosis involves fourth moments, it is sensitive to outliers: a single
extreme data point can dominate the estimate.
More robust contrast functions such as $G(x) = \log\cosh(x)$ (used in FastICA) are
preferred in practice when outliers are expected.
\end{notebox}


% ==========================================================
\chapter{Factor Analysis}
% ==========================================================

\section{The Factor Analysis Model}

% TODO

\section{Maximum Likelihood Factor Analysis}

% TODO

\section{Factor Analysis vs.\ PCA and ICA}

% TODO

\section{Artificial Factor Analysis Examples}

% TODO

\section{Real World Factor Analysis Examples}

\subsection{Iris Data Set}

% TODO

\subsection{Multiple Tissue Data Set}

% TODO


% ==========================================================
\chapter{Scaling and Projection Methods}
% ==========================================================

\section{Projection Pursuit}

% TODO

\section{Multidimensional Scaling}

\subsection{The Method}

% TODO

\subsection{Examples}

% TODO

\section{Non-negative Matrix Factorization}

\subsection{The Method}

% TODO

\subsection{Examples}

% TODO

\section{Locally Linear Embedding}

\subsection{The Method}

% TODO

\subsection{Examples}

% TODO

\section{Isomap}

\subsection{The Method}

% TODO

\subsection{Examples}

% TODO

\section{The Generative Topographic Mapping}

\subsection{The Method}

% TODO

\subsection{Examples}

% TODO

\section{\textit{t}-Distributed Stochastic Neighbor Embedding}

\subsection{The Method}

% TODO

\subsection{Examples}

% TODO

\section{Self-Organizing Maps}

\subsection{The Method}

% TODO

\subsection{Examples}

% TODO


% ==========================================================
\chapter{Clustering}
% ==========================================================

\section{Mixture Models}

\subsection{The Method}

% TODO

\subsection{Mixture of Gaussians}

% TODO

\subsection{Mixture of Poissons}

% TODO

\subsection{Examples}

% TODO

\section{\textit{k}-Means Clustering}

\subsection{The Method}

% TODO

\subsection{Examples}

% TODO

\section{Hierarchical Clustering}

\subsection{The Method}

% TODO

\subsection{Examples}

% TODO

\section{Similarity-Based Clustering}

\subsection{Aspect Model}

% TODO

\subsection{Affinity Propagation}

\subsubsection{The Method}

% TODO

\subsubsection{Examples}

% TODO

\subsection{Similarity-based Mixture Models}

\subsubsection{Similarity-based Mixture of Gaussians}

% TODO

\subsubsection{Representing the Centers}

% TODO

\subsubsection{Update Rules}

% TODO

\subsubsection{Examples}

% TODO


% ==========================================================
\chapter{Biclustering}
% ==========================================================

\section{Types of Biclusters}

% TODO

\section{Overview of Biclustering Methods}

% TODO

\section{FABIA Biclustering}

% TODO

\section{Examples}

% TODO


% ==========================================================
\chapter{Hidden Markov Models}
% ==========================================================

\section{Hidden Markov Models in Bioinformatics}

% TODO

\section{Hidden Markov Model Basics}

% TODO

\section{Expectation Maximization for HMM: Baum-Welch Algorithm}

% TODO

\section{Viterbi Algorithm}

% TODO

\section{Input Output Hidden Markov Models}

% TODO

\section{Factorial Hidden Markov Models}

% TODO

\section{Memory Input Output Factorial Hidden Markov Models}

% TODO

\section{Tricks of the Trade}

% TODO

\section{Profile Hidden Markov Models}

% TODO


% ==========================================================
\chapter{Boltzmann Machines}
% ==========================================================

\section{The Boltzmann Machine}

% TODO

\section{Learning in the Boltzmann Machine}

% TODO

\section{The Restricted Boltzmann Machine}

% TODO


\end{document}
